\section{Count-prefix growth and an all-depth stationary-resource separator order}
This section proves a new count-prefix inequality and an unbounded
order-transport reduction. It does NOT prove that the resulting
separator is nonnegative in every required case. General C, A2 and
ORIGINAL coefficient unimodality remain open.

\begin{theorem}[A Fibonacci growth law through the non-one count]
Let $F$ be generated from an initial $1^\alpha$, $\alpha\ge1$, by
$h$ capacities at least two. Their order need not be sorted. Then
\[
 F_n\ge F_{n-1}+F_{n-2}\qquad(2\le n\le h+1),
 \qquad 2F_n\ge3F_{n-1}\qquad(2\le n\le h+1).       \tag{1}
\]
The second statement uses the established actual/formal source
growth bound $F_{n-1}\le2F_{n-2}$. It is sharp, for example at
rank three on sufficiently long all-binary arms.
\end{theorem}
\begin{proof}
Write $S_t=z+\cdots+z^t$, $B=1-z-z^2$, and
$T_h=BF_h-1$. Initially
\[
 T_0=-z^{\alpha+1}-z^2[\alpha+1],
\]
so its coefficients below rank two vanish. For a capacity $t\ge2$,
\[
 T_{h+1}=(S_t-S_2)+S_tT_h.                         \tag{2}
\]
The first summand is nonnegative and the second raises the earliest
possible negative rank by at least one. Strongly, induction gives
$[z^n]T_h\ge0$ for every $0\le n\le h+1$.
Identifying the coefficient of $BF_h$ proves the first inequality;
the doubling bound proves the second. This is an unbounded word
induction, not extrapolation from finite counts.
\end{proof}

Let $d_n=F_n-F_{n-1}$ use zero extension, including $d_0=1,d_1=0$.
For the quantitative nonrecovery separator $Q=3M-e$, the SOURCE
variation weight is
\[
 W_n=3d_n-F_n.
\]
Equation (1) gives $W_n\ge0$ at $2\le n\le h+1$.
The two exceptional initial values are $W_0=2$, $W_1=-1$, while
$W_2=1$. Consequently EVERY uniform window of width at least two,
whose upper index is at most $h+1$, has nonnegative W sum.
If the window contains index one, it contains either index zero
(when it ends at one) or index two; those terms pay its sole negative
unit. Every other included term is nonnegative.

\begin{theorem}[Future source order for a quantitative separator]
Compare two predecessor source prefixes with equal leading jet
through normalized index $m$ and componentwise ordered later source
coefficients. Apply the SAME legal append $t\ge2$. Let the current
child first weight be $p$, opposite arm count $g$, and comparison
rank $N$. If
\[
                         N-p-m-1\le g+1,           \tag{3}
\]
then the child's $Q=3M-e$ cost increases with the predecessor source.
The strict positive-D prefix guard transports by native source order.
Current Q queries must be made for every terminal append to the
predecessor frontier; source pruning of the children themselves need
not preserve Q, since its current source weights can have a negative
coefficient at opposite rank one.
\end{theorem}
\begin{proof}
For a current source change the first variation of Q is W on the
opposite ordinary prefix. In the growth phase the predecessor source
variation is convolved with $[t]$; in the stationary phase it is
convolved with $S_t$. Thus every varied predecessor coefficient is
weighted by a uniform W window of width $t\ge2$. The largest possible
upper read is $N-p-m-1$ (stationary reads are one lower).
Equation (3) puts every such window in the preceding positive range.
The birth term agrees for the two compared states and cancels.
Apply coefficientwise source order and the native guard theorem.
\end{proof}

\subsection{The pre-count resource allocates the fixed jets at every depth}
At $N=p+q+s$ let the counts be $h,g$, with $N<h+g$.
The exact numbers of stationary appends are
\[
 \sigma_X=h-p+2,\qquad \sigma_Y=g-q+2,
 \qquad\sigma_X+\sigma_Y\ge s+5.                   \tag{4}
\]
First weight starts at two and every nonstationary legal append raises
it by one, proving these identities. Set
\[
 m_X=s-\min(s,\sigma_Y),\qquad
 m_Y=\min(s,\sigma_Y).
\]
Then $m_X+m_Y=s$, $m_X\ge\max(0,s-\sigma_Y)$ and
$m_Y\ge\max(0,s-\sigma_X)$. In particular
\[
 q+s-m_X-1\le g+1,\qquad
 p+s-m_Y-1\le h+1.                                \tag{5}
\]
Fix the source jets through $m_X,m_Y$. All unfixed cross terms in
M, e and Q vanish: their two normalized indices sum to at least
$m_X+m_Y+2=s+2>s$. The two residual Q costs are separable; (5)
places BOTH in the positive-parent window theorem.

Thus, at EVERY pre-count overlap depth, there is an explicit
count-dependent choice of fixed jets for which Q costs have a valid
native-source FUTURE simulation and preserved query guards.
This proves an order reduction and a verification scheme, NOT that
the Q lower table is positive. Negative-current profiles can have
negative Q, so a uniform positive Q assertion without a sign or
other branch condition is not supplied.

For a fixed queried N, counts can be capped at N while preserving
the pre-count resource: if a count was larger than N, replacing it
by N still leaves total count greater than N. Together with finite
source-prefix clipping and stationary forgetting, the scheme has
finite closure at every fixed depth/first-weight pair. The presently
implemented native grammar fixes TWO leading coefficients and caps
counts at six; it does not yet implement these variable jet lengths
or the N-count cap. This theorem is therefore an analytic next-step
reduction, not a claim that such general finite tables have been run.

\paragraph{A tempting stronger binary lower bound is false.}
Finite small-rank minima equal Fibonacci numbers, but they cannot be
extended as a theorem $w_n\ge\mathrm{Fib}_{n-4}$ through every binary
FIRST D mode. The ACTUAL sorted arm $1^{120}2^{40}$ has first D mode
59 and there
\[
 w_{59}=104121337030<139583862445=\mathrm{Fib}_{55}.
\]
Its current derivative strictly increases. This rejects only that
stronger proposed lower bound. The completed binary fourth-overlap
proof uses a complete grid through rank eighteen and does not use the
false unbounded extrapolation. General UI/C/A2 are not refuted.

\paragraph{Scoped formalization.}
\texttt{CountPrefixSeparator.lean} verifies six unbounded scoped
kernels: the complete two-variable positive-cascade induction,
the doubling-to-three-halves implication, arbitrary finite-window
single-negative compensation, interval geometry, explicit stationary
budget allocation, and both count-prefix read bounds. Source/word
coefficient identification, formal-series variation and the full
separator-order composition are the written proofs above.
